\chapter[Perception: World or Interpretation?]{Do We See the World, or the Brain's Interpretation?}
\section{A ``Bent'' Chopstick}
Pick up an ordinary chopstick.

Place it obliquely in half a glass of water and look down. It breaks at the surface, the submerged and exposed portions meeting at an angle, as if bent and awkwardly rejoined. Remove it: straight and intact. Replace it: bent again.

You know it has not bent. That is not the point. Your \textbf{eyes} insist it has. Physics knowledge and confidence cannot remove the visible kink. You see a bent stick and believe in a straight one: seeing and believing have separated inside your head.

Try other free experiments. Hold the book at arm's length, fix one eye on a nearby point and move a printed word slowly into peripheral vision. At one position it disappears. Each eye has a blind spot without photoreceptors, yet your visual field has never seemed punctured. Before switching on lights at dusk, white paper takes orange-red sunset light but still looks white. An opening door no longer projects a rectangle on your retina, yet remains a rectangular door. Psychologists call this \textbf{perceptual constancy}: changing light and angles do not unsettle perception's stable world.

Convenient, and suspicious. Without stability, every opening door's changing shape would drive you mad. Yet perception clearly does more than copy: it processes, modifies and concludes without permission.

Such traces abound. Equal lines with differently directed arrowheads look unequal, even after a ruler demonstrates equality. Films contain no actual screen movement, only twenty-odd still frames per second, yet vision supplies continuous running, falling and embracing. The moon illusion makes a horizon Moon look far larger than an overhead one, while same-night photographs show equal sizes. Eyes, rulers, cameras and physics hold daily meetings of conflicting reports.

Is the world before you---table, water's colour, tree outside---its actual appearance, or an interpretative draft supplied by your brain?

\section{The Lesson That Erupted}
Come into a scene.

Early summer, the first afternoon lesson in a secondary classroom. Half-drawn curtains meet a projector beam carrying slowly drifting chalk dust. The board bears Seeing is believing and a large question mark.

The biology teacher projects a chequerboard with a green cylinder casting a shadow. Pointing to dark-looking square A outside and light-looking square B inside, the teacher asks which is darker.

Everyone answers A. Lin Yuan, in the third row, shouts loudest: he believes pictures establish truth.

Without replying, the teacher samples A and B with a colour picker and places their patches side by side at the screen's centre.

Silence for two seconds, then uproar.

The colours are identical.

Impossible. The instrument must be broken. The shadowed square is obviously white. Lin Yuan stands, demanding a paper strip connecting them. The teacher obliges: its colour runs uniformly from A to B without a seam. Remove it and A darkens again while B lightens.

By the window, Xu Tang speaks softly. In painting lessons, she learned to lighten a colour placed in a dark area. Perhaps eyes are not merely deceiving, but constantly correcting.

Lin Yuan objects: is that not deception? The instrument measures objective colour while eyes show falsehood. What good are eyes?

The bell rings; nobody moves. The teacher leaves one instruction: show your parents tonight, then ask whether your eyes have malfunctioned or are doing something else.

Remember this room. Our chapter addresses Lin Yuan's anger and the teacher's question.

\section{Can Eyes Be Trusted?}
Set out the conflict before judging.

Lin Yuan's position is ancient and respectable: \textbf{seeing is believing}. The Book of Han attributes to veteran general Zhao Chongguo the saying that a hundred reports cannot equal one sight. Courts want physical evidence, journalists photographs, shoppers something to weigh in a hand. Direct observation as a gold standard underpins modern science: bacteria through microscopes and Jupiter's moons through telescopes. If eyes are untrustworthy, what remains?

Opposite stands difficult evidence: submerged chopsticks, chequerboard squares, desert mirages and dream corridors whose texture seems touchable. Dreams likewise show and persuade, until waking reveals a film with one viewer also directing, filming and playing every role. Your perceptual system can manufacture a convincing world from nothing. That is not merely philosophical speculation; you experienced it last night.

Dreams deserve attention because they are easily dismissed. Almost everyone dreams, even without morning recall. Dreams have sights, sounds, feelings and coherent plots. Some dreamers pinch themselves and solemnly conclude they are awake, while still dreaming. Even the feeling of verification can be fabricated offline. Does a production line capable of fooling humanity at night really close by day?

The deeper question emerges:

\textbf{Do we see the world itself, or the brain's interpretation?}

If the former, explain illusions, dreams and mirages. If the latter, how do individual brain-produced drafts place us in one shared world? Why is science true rather than a collective high-definition dream?

Do not choose yet. Present both cases fully first.

\section{Whose Reasons Are Stronger?}
\textbf{First: seeing is believing.} Give this often-dismissed-as-naive position a fair hearing.

First, perception is usually accurate. A chopstick appears bent, yet your grasp reaches its real position. Doors change appearance, yet you avoid their frames. A constantly mistaken system could not navigate rush-hour traffic or let surgeons stitch millimetre-wide vessels. Second, senses cross-check. Eyes report bending, hands and rulers report straightness; two witnesses against one settle it. Illusions need not be dead ends because the world permits comparison. Third, science rests on shared observation: different people see bacteria and moons. Purely private perception would make reproducible experiments impossible.

Strong reasons, but the opposition is strong too.

\textbf{Second: eyes are advocates, not witnesses.} Advocates present clients' best cases, not necessarily truth. Perception's client is survival. Knowing fruit ripeness and approaching tigers matters more than wavelengths. Speed, energy economy and survival encourage shortcuts: filling blind spots, correcting shadowed colours and turning peripheral darkness into a snake. The classroom squares expose unauthorised interpretation. Dreams are worse: not a blade of grass in last night's corridor came from external input. Your brain can render complete virtual reality without your noticing its operation.

\textbf{Third: practical trust.} Pilots, emergency physicians and referees offer a neglected position. Knowing eyes and instruments fail, they must act immediately. Their answer is procedure, not allegiance: cross-checks, two-person verification and later remeasurement. Doubt belongs in review, trust in immediate action.

Consider another tradition. Master Lü's Spring and Autumn Annals tells how Confucius's party ran short of food between Chen and Cai. Yan Hui obtained rice and cooked it. Confucius saw him take some from the pot and eat. Saying nothing until it was ready, he proposed offering clean rice to ancestors; rice handled this way was ritually unsuitable. Yan Hui explained that soot had fallen in and he ate the contaminated portion rather than waste it. Confucius remarked, broadly, that even trusted eyes could mislead.

What failed? Not his eyes: Yan Hui did eat rice. The \textbf{interpretation} turned eating into stealing. Eyes deliver; brains write, and drafts often err. Language offers another example. Japanese ao covers blue sky, green signals and unripe fruit. Ao shingō, a green traffic light, does not imply colour-blindness: the linguistic boundary differs from ours. A seemingly natural colour division moves between languages.

\section[Thinking Bridge: Three Viewing Layers]{Thinking Bridge: Live Views, Relays and Eager Commentators}
Turn this into a portable tool. When someone insists their own eyes cannot be wrong, separate perception into three layers before judging:

\textbf{First: the world itself.} A and B supply physically equal grey levels. This is not negotiable.

\textbf{Second: sensory signals.} Light reaches the retina and becomes electrical activity, already incomplete. Blind spots exist, peripheral colour discrimination is weak, and eyes receive only a narrow electromagnetic band. Bees see ultraviolet that you cannot: the signal exists, but your receiver lacks its frequency range.

\textbf{Third: the brain's interpretation.} Experience, context and expectations shape the final seen world: shadowed squares brighten, blind spots fill and half-open doors regain rectangularity.

These layers make three philosophical positions into understandable viewing modes:

\begin{itemize}
\item \textbf{Direct realism} treats perception as a \textbf{live broadcast}. We see the world itself without an intervening draft: squares' colours and tables' hardness present themselves directly. Intuitive and supportive of common sense, it struggles with dreams and illusions: how does a live signal acquire nonexistent content?
\item \textbf{Representationalism} treats perception as a \textbf{relay}. What we directly see is an internally generated representation, with the world behind it. Like viewing a security monitor, we infer a corridor outside. This explains erroneous pictures, but asks how permanent screen-viewers can check the corridor independently.
\item \textbf{Predictive perception}, influential in recent cognitive science, resembles an \textbf{over-eager commentator}. Rather than wait passively, the brain \textbf{predicts} from experience and uses incoming signals to correct errors. Correct predictions need little upward reporting, explaining unnoticed familiar routes. Errors rush upwards and trigger revisions. The shadowed square elicits a premature brightening prediction. Dreams are commentary broadcasting alone, rendering stored experience without input.
\end{itemize}
Apply the tool to a mirage. Water and reflections seem to appear on desert ground or a hot road. First, no water exists: hot ground heats air, and light bends through temperature layers. Second, eyes accurately receive bent light from the sky at a low angle. Third, experience associates low bright patches with water reflections, prompting the draft to announce water. Neither desert nor eyes deceive; a usually useful expectation loses its race against uncooperative physics.

Apply it to people too. For an eyewitness claim, ask what objectively happened, how light, angle, distance or equipment limited signals, and how expectations, commitments and experience interpreted them. Many apparent disputes over observation concern the third-layer draft, not the first-layer event.

\section{Butterflies, Demons and Vats}
Read primary material. Humans have seriously questioned perception in writing for over two millennia.

Zhuang Zhou's famous Warring States dream appears in Making All Things Equal:

\begin{quote}
Once Zhuang Zhou dreamed he was a butterfly, fluttering freely, content in its wishes, knowing nothing of Zhou. Suddenly he awoke, unmistakably Zhou. Was Zhou dreaming himself a butterfly, or was a butterfly dreaming itself Zhou?
\end{quote}
In brief, Zhou dreams happy butterfly flight without knowing himself, then wakes as Zhou lying stiffly in bed. Which is dreaming which?

Notice he supplies no answer. The passage concludes that Zhou and butterfly must differ, calling this the transformation of things. It is not wordplay: dreamlike freedom and waking recognition carry \textbf{the same feeling of reality}. You do not know you are dreaming within a dream; otherwise it would not be one. Since reality-feeling occurs in unreal scenes, it cannot judge truth.

Plato offers a structurally related allegory in Republic VII, paraphrased here. Prisoners chained from childhood see only wall shadows and regard them as all reality. One escapes into sunlight and real things, then returns to disbelief. Shadows relate to objects as perception relates to world: relayed images, differing in whether you can change cameras.

A millennium and more later, Descartes radicalises doubt in Meditations. Might a powerful, cunning demon supply the illusion of sky, earth and body? His purpose is a stress test, doubting everything possible until something survives. Deception itself establishes a thinking entity being deceived, the foundation behind I think, therefore I am.

In the twentieth century, Hilary Putnam gives the demon technological equipment in the \textbf{brain-in-a-vat} experiment. An extracted brain floats in nutrients, connected to a supercomputer supplying all normal signals. It sees furniture, touches a cat and experiences a seamless lifetime. How can you prove you are not that brain?

This is not uniquely Western obsession. Earlier Buddhist thought offers related formulations, including the Diamond Sutra's closing verse:

\begin{quote}
All conditioned things are like dreams, illusions, bubbles and shadows, like dew and lightning. Contemplate them in this way.
\end{quote}
Condition-dependent things are fleeting and insubstantial. Distinguish purposes, however: Buddhism discourages attachment, Descartes seeks foundations, Putnam investigates language and meaning. Similar doubt serves three different projects. Thought experiments matter through their use, not their frightening quality.

\section{Why Illusions Occur: Three Accounts}
Return to the classroom. Why do A and B look different? Three genuinely different explanations address the same fact.

\textbf{First: faulty equipment.} Eyes are precise but imperfect cameras, with evolutionary bugs not yet repaired. Simple and intuitive, this explains some phenomena, including the blind spot's wiring defect. But consider a \textbf{counterexample}: colour sampling and a connecting strip establish equality, yet removing the strip restores the illusion. Calibrated faulty equipment should recover, but vision \textbf{refuses calibration}. It seems committed to an algorithm rather than broken. Knowing does not cure mis-seeing, suggesting algorithmic by-product instead of equipment error.

\textbf{Second: optimal guessing, or predictive perception.} The brain supplies useful guesses quickly, not exact copies. Natural shadows reduce reflected light without changing surface colour, prompting automatic compensation. The chequerboard activates that rule despite being an artificial laboratory oddity. Illusion is a well-tested survival algorithm out of context, not a defect. This predicts when and in which direction error occurs: overcorrection rather than randomness. But does the brain literally compute predictions, or is computer-like language simply our handiest metaphor? That remains unsettled.

\textbf{Third: good enough, or evolutionary usefulness.} Some cognitive scientists propose that evolution values usefulness, not truth. Perception resembles desktop icons: a blue square representing a file need not resemble the file itself. Colours, shapes and even space may be interfaces related to reality without matching it. This can encompass constancy, illusions and culturally differing colour categories. But if everything is an icon, are instruments merely another icon layer? An account explaining so much invites suspicion about what it discards. Icons must connect to files; usefulness and resemblance may be harder to separate than claimed.

Each account holds part, not all, of the truth. Even explaining interpretation generates competing drafts of the same world.

\section[Further Challenge: Living Scepticism]{Further Challenge: Can Scepticism Be Lived Consistently?}
Dreams, demons and vats together pose \textbf{scepticism}, systematic doubt about knowing the world. Intellectually exhilarating, it meets a simple question: \textbf{can it be practised?}

Not merely conceived, but lived.

Pyrrho reportedly suspended every judgement. Diogenes Laertius's Lives of Eminent Philosophers, written centuries later, depicts friends protecting him from carts and cliffs because he refused even to believe an approaching cart. Discount the story's authenticity; it is probably an opponent's caricature. Yet it exposes a weakness: \textbf{theory can doubt the road; the body must cross it now}. A philosophy impossible to practise in those four seconds is incomplete or insincere, its advocate not really believing it.

Hume saw deeper. Radical doubts may survive in a study, but dining, playing and talking with friends dissolve them within an hour. Life disarms rather than refutes them. Scepticism neither wins nor loses; its proper effect is cooling \textbf{overconfident} beliefs, not ending belief.

G. E. Moore offered another answer. Asked to prove an external world, he raised one hand, then the other, and concluded the external world exists. Apparently simplistic, his point was comparative certainty: sceptical premises are less secure than having two hands. Using weaker premises to overturn stronger convictions is a poor bargain. Doubt needs reasons no more doubtful than its targets.

Where did Descartes's demon lose? Perhaps nowhere; it is merely \textbf{unusable}. Total doubt resembles full-power radar, incompatible with shopping, examinations and friendship. That does not make scepticism worthless. \textbf{Local, specifically addressed doubt} is exceptionally powerful. Science doubts particular conclusions under rules, changing when better evidence arrives. Descartes's lasting legacy is less demon than procedure: examine beliefs individually, retaining justified ones and downgrading others.

\section{When Instruments Become New Senses}
The old question now rests on your nose and in your pocket.

First, extension. Microscopes take eyes into worlds a thousandth of a hair's diameter. Seventeenth-century Dutch observer Leeuwenhoek saw aquatic microorganisms with homemade lenses, calling them animalcules. His reports to the Royal Society met disbelief, repeated checking and independent observation before a drop's city of life became accepted. Telescopes reached Moon and Jupiter; Galileo's orbiting Jovian moons undermined the claim that everything circles Earth. Infrared thermometers, CT scanners, seismometers and gravitational-wave detectors are new senses revealing what unaided eyes cannot in principle see.

Instruments even redefine normal senses. Corrective glasses make a classroom board clear, without anyone calling their mediated world false. Hearing aids convert sound into signals; cochlear implants bypass the ear to stimulate nerves directly. Reliability depends not on mediation's absence, but on its stability and agreement with other evidence.

But \textbf{instruments do not merely extend senses; they reinterpret them}. Galileo's difficulty was whether telescopic sights counted, not simply whether he could see. Were they reality or instrument-made phantoms? The objection was sensible. The answer required procedures, not infallibility: agreement across instruments, confirmation by different principles such as cultured bacteria matching microscope observations, and successful predictions. Science upgraded observation by inventing \textbf{institutions through which different eyes scrutinise one another}, not flawless eyes. Instruments yield interpretations subjected to repeated questioning, much like brains, but with the questioning exposed publicly.

In daily life, default camera beautification already supplies an algorithmic face. Video recommendations render a personal world from watch time and likes, predicting preferences and cutting reality accordingly. You often see personalised commentary rather than the world. Deepfakes remove pictures' final guarantee: high-definition video alone no longer suffices as evidence.

This does not justify distrusting everything because everything is interpreted. That is lazy and dangerous. Colour-picker checks produce a \textbf{defensible} equality between A and B. Persistent flat-Earth belief does not become equivalent to a spherical-Earth account. Interpretations differ in quality through cross-checking and predictive success. Brains have survived by such standards for millions of years; science built the modern world with them. Apply them to algorithm-fed worlds too.

\section[Your Turn: Remove the Glasses?]{Your Turn: Should You Take Off Those Glasses?}
Return to the classroom.

Before leaving, Lin Yuan asks: if eyes correct, instruments interpret and algorithms personalise, is what we see still the world?

The teacher does not answer. Now you must, with reasons. First weigh the chapter again.

On one side lie Zhou's waking bewilderment, Descartes's demon and the vatted brain. They show the absence of an absolute authenticity certificate proving this is not a high-definition dream.

On the other stand traffic demanding immediate trust, microscopes answering culture dishes and equal colour samples. Interpretations really differ, and some withstand worldwide cross-examination.

Between them lies perceptual constancy, correcting the world so you recognise shadowed faces but misjudge chequerboards: one algorithm, half bodyguard and half deceiver.

Should seeing is believing be discarded or revised? How would you revise it? Further, would you spend your life in flawless virtual reality offering any desired existence? If not, what could you not surrender: reality itself, or something you cannot yet articulate?

There is no standard answer. Your thought now is another interpretative draft from your brain. The difference is that, from today, it bears your own annotations.
